% theme: power_tangent_chord_and_incircle
% 方べき・接弦定理は発展事項として扱う。表：相似による方べき、接弦定理。裏：接線長、内接円半径。
\MajorQuestionLesson{方べきの定理と接弦定理・内接円}
\SetRowSpace{2mm}

\TypeQuestion{図から相似な三角形を見つけ、$x$ の値を求めなさい。長さの単位は $\mathrm{cm}$ とする。}{power_of_point_by_similarity}
\QQFigRowThree[31mm]{fig/lesson_03_power_01}{}{fig/lesson_03_power_02}{}{fig/lesson_03_power_03}{}
\QQFigRowThree[31mm]{fig/lesson_03_power_04}{}{fig/lesson_03_power_05}{}{fig/lesson_03_power_06}{}

\TypeQuestion{図の直線は円の接線である。接弦定理を用いて、$x$ の大きさを求めなさい。}{tangent_chord_theorem}
\QQFigRowThree[31mm]{fig/lesson_03_tangent_01}{}{fig/lesson_03_tangent_02}{}{fig/lesson_03_tangent_03}{}
\QQFigRowThree[31mm]{fig/lesson_03_tangent_04}{}{fig/lesson_03_tangent_05}{}{fig/lesson_03_tangent_06}{}

\AnswerColumnBreak
\LessonQuestionBreak

\TypeQuestion{円は図形の各辺に接している。同じ点から引いた接線の長さが等しいことを用いて、$x$ の値を求めなさい。長さの単位は $\mathrm{cm}$ とする。}{equal_tangent_lengths}
\QQFigRowThree[31mm]{fig/lesson_03_incircle_len_01}{}{fig/lesson_03_incircle_len_02}{}{fig/lesson_03_incircle_len_03}{}
\QQFigRowThree[31mm]{fig/lesson_03_incircle_len_04}{}{fig/lesson_03_incircle_len_05}{}{fig/lesson_03_incircle_len_06}{}

\TypeQuestion{三角形の3辺の長さと面積を用いて、内接円の半径 $r$ を求めなさい。長さの単位は $\mathrm{cm}$、面積の単位は $\mathrm{cm}^2$ とする。}{incircle_radius_from_area}
\QQFigRowThree[31mm]{fig/lesson_03_incircle_r_01}{}{fig/lesson_03_incircle_r_02}{}{fig/lesson_03_incircle_r_03}{}
\QQFigRowThree[31mm]{fig/lesson_03_incircle_r_04}{}{fig/lesson_03_incircle_r_05}{}{fig/lesson_03_incircle_r_06}{}
